\documentclass[11pt]{article}
\usepackage[margin=25mm]{geometry}
\usepackage{fontspec}
\setmainfont{Latin Modern Roman}
\usepackage{amsmath,amssymb,mathtools}
\usepackage{unicode-math}
\setmathfont{Latin Modern Math}
\usepackage{xurl,hyperref}
\hypersetup{colorlinks=true,urlcolor=blue}
\setcounter{secnumdepth}{0}
\setlength{\parindent}{0pt}
\setlength{\parskip}{5pt}
\setlength{\emergencystretch}{3em}
\newcommand{\sref}[2]{\hyperlink{source:#1:#2}{[S#2]}}
\newcommand{\cataloguescope}[1]{\paragraph{Recorded target.}#1}
\newcommand{\R}{\mathbb R}
\newcommand{\C}{\mathbb C}
\newcommand{\N}{\mathbb N}
\newcommand{\Z}{\mathbb Z}
\newcommand{\Q}{\mathbb Q}
\newcommand{\F}{\mathbb F}
\newcommand{\eps}{\varepsilon}
\begin{document}
% UnsolvedMath ID 20002527; source catalogue code AIM-PHYSICS-0035
\setcounter{section}{20002526}
\section{The L2-signature of toric hyperkahler quotients}\label{20002527}
{\small UnsolvedMath ID 20002527 · Mathematical Physics}
\subsection{Definitions and mathematical statement}

\paragraph{Shared notation.}$\mathbb N=\{1,2,\ldots\}$; the standard number fields have their usual meanings. All additional conventions are stated locally.


\paragraph{Definitions and notation.}A hyper-Kähler metric has three parallel complex structures satisfying the quaternionic relations. In four dimensions a gravitational instanton is a complete noncompact Ricci-flat geometry in the source’s self-dual setting. ALE and ALF describe different asymptotic geometries; ALF has an asymptotic circle direction over a three-dimensional base. For a Riemannian manifold $(M,g)$, $L^2$ norms use the metric volume and pointwise metric on forms. Reduced $L^2$ cohomology is $\ker d/\overline{\operatorname{im}d}$ for the specified closed $L^2$ complex. On a complete manifold its harmonic representatives satisfy $d\alpha=d^*\alpha=0$ and $\int_M|\alpha|^2<\infty$. In even middle degree, the Hodge star splits harmonic forms into self-dual and anti-self-dual sectors. Their dimension difference is an $L^2$ signature when finite and suitably defined. An index theorem must specify the Fredholm operator and the asymptotic or boundary corrections needed on a noncompact space. Toric hyper-Kähler quotients have combinatorial data associated with the reducing torus.

\paragraph{Catalogue wording retained for traceability.}
For the source’s toric hyper-Kähler quotients, understand the $L^2$ signature and its sign, using the source’s middle-degree harmonic-form context and its relation to combinatorics and matroids. More generally, prove an appropriate index theorem for these manifolds and determine whether contributions from subsystems occur.
\paragraph{Scope and source qualification.}The assertions about concentration in middle degree and self-duality are historical source context with geometric assumptions, not newly verified universal statements. This is a third first step in the shared programme of 20002525–20002527. The catalogue title/status was paired with a machine-generated research attempt; that attempt is not peer-reviewed evidence of settlement. The statement covers the original source questions rather than only the title’s subcase.
\subsection{Short English statement}
Compute L² signatures and suitable index formulae for toric hyper-Kähler quotients, including possible subsystem terms.
\subsection{Research attempt}
\paragraph{Model, date and investigated scope.}
OpenAI GPT-6 Astra (configured model \texttt{gpt-6-astra}), ultra reasoning; generated 11 October 2026. Proved the conditional L2 intersection-form sign decomposition and exposed boundary terms in cohomological replacement. No toric quotient signature, matroid formula or subsystem index theorem is obtained.

\paragraph{Sources and literature check.}
Read the supplied complete definition, catalogue wording, scope qualifications and source list for this canonical ID. The original-source citations are inherited as supplied references; no fresh primary-source or novelty search was performed for this attempt. No current open-status claim is inferred from this calculation.

\paragraph{A sign criterion without an index theorem.}Let \(M\) be an oriented complete Riemannian manifold of dimension \(4m\), and suppose the real space \(\mathcal H\) of square-integrable harmonic \(2m\)-forms is finite dimensional. The Hodge star satisfies \(*^2=1\) in this degree, is an isometry and preserves harmonic forms, so \(\mathcal H=\mathcal H_+\oplus\mathcal H_-\). The integral pairing is
\[Q(\alpha,\beta)=\int_M\alpha\wedge\beta=\langle\alpha,*\beta\rangle_{L^2}.
\]
It converges absolutely by Cauchy--Schwarz. On the two eigenspaces it is respectively the positive and negative L2 inner product; cross terms vanish because distinct eigenspaces of a self-adjoint involution are orthogonal. Hence \(\operatorname{sign}Q=\dim\mathcal H_+-\dim\mathcal H_-\). If a geometric argument places all middle harmonic forms in one star eigenspace, the sign is forced, but that placement is an extra theorem.
\paragraph{Boundary control cannot be skipped.}On a compact truncation \(M_R\), replacing a closed form by \(\alpha+d\eta\) changes its pairing with a closed \(\beta\) by \(\int_{\partial M_R}\eta\wedge\beta\), by Stokes. Square integrability of \(\alpha,\beta\) alone does not specify an admissible primitive \(\eta\) or force that boundary term to vanish. Therefore importing a compact topological signature formula requires additional decay/range control.
\paragraph{Gap.}No finite-dimensionality, star concentration, matroid formula or Fredholm index theorem is proved for the source's toric quotients. Subsidiary cluster ends could contribute precisely through the missing boundary analysis. The first step still needed is a controlled domain/asymptotic description of the signature operator. The elementary sign identity only states what follows once the analytic geometry has been established.

\paragraph{Verification and adversarial audit.}
Checked dimension 4m star square, absolute pairing convergence, orthogonal eigenspaces, finite-dimensional assumption and Stokes boundary term. This was an LLM adversarial audit, not independent human verification; no theorem prover was used. Prompts A then B applied to this ID: Checked dimension 4m star square, absolute pairing convergence, orthogonal eigenspaces, finite-dimensional assumption and Stokes boundary term. Scope remains limited: Establish Fredholm/domain and end-decay control for the signature operator, including any cluster boundary contributions.

Coordinator checked the middle-degree star pairing, its signature decomposition and the Stokes boundary term needing decay control. This additional check was model-assisted, not independent human verification.

\paragraph{Tools, process and contributions.}
Python 3 and local shell used to read supplied files and serialize this manuscript; mathematics derived symbolically in the text, with no numerical experiment or formal proof assistant. Contribution prompt A was followed by prompt B. GPT-6 Astra with ultra reasoning drafted the research and performed its stated self-audit; any additional coordinator contribution is identified in the verification and audit. A GPT-6-based Codex coordinator assembled the manuscript and checked repository integration; its exact variant and reasoning configuration were not exposed. No human mathematical contribution, independent human proof check or formal proof certification is supplied. Complete derivations and the consulted source references are included here; private reasoning transcripts are not published.

\paragraph{Final status and first remaining gap.}
\textbf{UNRESOLVED}. Establish Fredholm/domain and end-decay control for the signature operator, including any cluster boundary contributions.

\section{Completion Estimate}
\noindent\textbf{COMPLETION: 0\%}
The calculation is an elementary benchmark or conditional reduction; it establishes no claimed new portion of the complete source target.

\subsection{Sources}
\begin{itemize}
\item[\hypertarget{source:20002527:1}{S1}] ULAM.AI and the UnsolvedMath Contributors, UnsolvedMath, record 20002527 (AIM-PHYSICS-0035). Catalogue curation under CC BY 4.0; definitions and English exposition added locally.
\url{https://huggingface.co/datasets/ulamai/UnsolvedMath}.
\item[\hypertarget{source:20002527:2}{S2}] American Institute of Mathematics, $L^2$ harmonic forms in geometry and string theory; notes by Anda Degeratu and Mark Haskins, March 2004, 21 March 2004, Hausel wrap-up programme and first step 3.
\url{https://aimath.org/WWN/l2harmonic/l2harmonic.pdf}.
\end{itemize}
\label{end:20002527}
\end{document}
